%!TEX program = pdflatex
%%%%%%%%%%%%%%%%%%%%%%%%%%%%%%%%%%%%%%%%%
% The Legrand Orange Book
% LaTeX Template
% Version 1.4 (12/4/14)
%
% This template has been downloaded from:
% http://www.LaTeXTemplates.com
%
% Original author:
% Mathias Legrand (legrand.mathias@gmail.com)
%
% License:
% CC BY-NC-SA 3.0 (http://creativecommons.org/licenses/by-nc-sa/3.0/)
%
% Compiling this template:
% This template uses biber for its bibliography and makeindex for its index.
% When you first open the template, compile it from the command line with the 
% commands below to make sure your LaTeX distribution is configured correctly:
%
% 1) pdflatex main
% 2) makeindex main.idx -s StyleInd.ist
% 3) biber main
% 4) pdflatex main x 2
%
% After this, when you wish to update the bibliography/index use the appropriate
% command above and make sure to compile with pdflatex several times 
% afterwards to propagate your changes to the document.
%
% This template also uses a number of packages which may need to be
% updated to the newest versions for the template to compile. It is strongly
% recommended you update your LaTeX distribution if you have any
% compilation errors.
%
% Important note:
% Chapter heading images should have a 2:1 width:height ratio,
% e.g. 920px width and 460px height.
%
%%%%%%%%%%%%%%%%%%%%%%%%%%%%%%%%%%%%%%%%%

%----------------------------------------------------------------------------------------
%	PACKAGES AND OTHER DOCUMENT CONFIGURATIONS
%----------------------------------------------------------------------------------------

\documentclass[11pt,fleqn]{book} % Default font size and left-justified equations

\usepackage[top=3cm,bottom=3cm,left=3.2cm,right=3.2cm,headsep=10pt,a4paper]{geometry} % Page margins

\usepackage{xcolor} % Required for specifying colors by name
\definecolor{blue}{rgb}{0.0, 0.18, 0.39}

% Font Settings
\usepackage{avant} % Use the Avantgarde font for headings
%\usepackage{times} % Use the Times font for headings
\usepackage{mathptmx} % Use the Adobe Times Roman as the default text font together with math symbols from the Sym­bol, Chancery and Com­puter Modern fonts

\usepackage{microtype} % Slightly tweak font spacing for aesthetics
\usepackage[utf8]{inputenc} % Required for including letters with accents
\usepackage[T1]{fontenc} % Use 8-bit encoding that has 256 glyphs
\hyphenation{Mi-nis-té-ri-o}


% Index
\usepackage{calc} % For simpler calculation - used for spacing the index letter headings correctly
\usepackage{makeidx} % Required to make an index
\makeindex % Tells LaTeX to create the files required for indexing
\usepackage{verbatim}

\usepackage[colorinlistoftodos,prependcaption,textsize=tiny,linecolor=red,backgroundcolor=red!25,bordercolor=red]{todonotes}
\usepackage{epigraph}
\renewcommand{\textflush}{flushepinormal}
\setlength{\epigraphwidth}{0.8\textwidth}

\usepackage{nameref}
\usepackage{booktabs}
\usepackage{graphicx}
\usepackage{float}
\usepackage{multirow}


% Bibliography
%\usepackage[backend=biber,style=authoryear,autocite=inline, citestyle=authoryear]{biblatex}
\usepackage[style=abnt]{biblatex}
\addbibresource{bibliography.bib} % BibTeX bibliography file
\defbibheading{bibempty}{}
\renewcommand*{\nameyeardelim}{\addcomma\space}

\newcommand{\VER}[1]{\begingroup\color{red}#1\endgroup}

%----------------------------------------------------------------------------------------

%----------------------------------------------------------------------------------------
%	VARIOUS REQUIRED PACKAGES
%----------------------------------------------------------------------------------------

\usepackage{titlesec} % Allows customization of titles

\usepackage{graphicx} % Required for including pictures
\graphicspath{{Pictures/}} % Specifies the directory where pictures are stored

\usepackage{lipsum} % Inserts dummy text

\usepackage{tikz} % Required for drawing custom shapes

\usepackage[english,brazil]{babel} % English language/hyphenation

\usepackage{enumitem} % Customize lists
\setlist{nolistsep} % Reduce spacing between bullet points and numbered lists

\usepackage{booktabs} % Required for nicer horizontal rules in tables

\usepackage{eso-pic} % Required for specifying an image background in the title page

%----------------------------------------------------------------------------------------
%	MAIN TABLE OF CONTENTS
%----------------------------------------------------------------------------------------

\usepackage{titletoc} % Required for manipulating the table of contents

\contentsmargin{0cm} % Removes the default margin
% Chapter text styling
\titlecontents{chapter}[1.25cm] % Indentation
{\addvspace{15pt}\large\sffamily\bfseries} % Spacing and font options for chapters
{\color{blue!60}\contentslabel[\Large\thecontentslabel]{1.25cm}\color{blue}} % Chapter number
{}  
{\color{blue!60}\normalsize\sffamily\bfseries\;\titlerule*[.5pc]{.}\;\thecontentspage} % Page number
% Section text styling
\titlecontents{section}[1.25cm] % Indentation
{\addvspace{5pt}\sffamily\bfseries} % Spacing and font options for sections
{\contentslabel[\thecontentslabel]{1.25cm}} % Section number
{}
{\sffamily\hfill\color{black}\thecontentspage} % Page number
[]
% Subsection text styling
\titlecontents{subsection}[1.25cm] % Indentation
{\addvspace{15pt}\sffamily\small} % Spacing and font options for subsections
{\contentslabel[\thecontentslabel]{1.25cm}} % Subsection number
{}
{\sffamily\;\titlerule*[.5pc]{.}\;\thecontentspage} % Page number
[] 

%----------------------------------------------------------------------------------------
%	MINI TABLE OF CONTENTS IN CHAPTER HEADS
%----------------------------------------------------------------------------------------

% Section text styling
\titlecontents{lsection}[.05em] % Indendating
{\footnotesize\sffamily} % Font settings
{}
{}
{}

% Subsection text styling
\titlecontents{lsubsection}[.5em] % Indentation
{\normalfont\footnotesize\sffamily} % Font settings
{}
{}
{}
 
%----------------------------------------------------------------------------------------
%	PAGE HEADERS
%----------------------------------------------------------------------------------------

\usepackage{fancyhdr} % Required for header and footer configuration

\pagestyle{fancy}
\renewcommand{\chaptermark}[1]{\markboth{\sffamily\normalsize\bfseries\chaptername\ \thechapter.\ #1}{}} % Chapter text font settings
\renewcommand{\sectionmark}[1]{\markright{\sffamily\normalsize\thesection\hspace{10pt}#1}{}} % Section text font settings
\fancyhf{} \fancyhead[LE,RO]{\sffamily\normalsize\thepage} % Font setting for the page number in the header
\fancyhead[LO]{\rightmark} % Print the nearest section name on the left side of odd pages
\fancyhead[RE]{\leftmark} % Print the current chapter name on the right side of even pages
\renewcommand{\headrulewidth}{0.5pt} % Width of the rule under the header
\addtolength{\headheight}{2.5pt} % Increase the spacing around the header slightly
\renewcommand{\footrulewidth}{0pt} % Removes the rule in the footer
\fancypagestyle{plain}{\fancyhead{}\renewcommand{\headrulewidth}{0pt}} % Style for when a plain pagestyle is specified

% Removes the header from odd empty pages at the end of chapters
\makeatletter
\renewcommand{\cleardoublepage}{
\clearpage\ifodd\c@page\else
\hbox{}
\vspace*{\fill}
\thispagestyle{empty}
\newpage
\fi}

%----------------------------------------------------------------------------------------
%	THEOREM STYLES
%----------------------------------------------------------------------------------------

\usepackage{amsmath,amsfonts,amssymb,amsthm} % For math equations, theorems, symbols, etc

\newcommand{\intoo}[2]{\mathopen{]}#1\,;#2\mathclose{[}}
\newcommand{\ud}{\mathop{\mathrm{{}d}}\mathopen{}}
\newcommand{\intff}[2]{\mathopen{[}#1\,;#2\mathclose{]}}
\newtheorem{notation}{Notation}[chapter]

%%%%%%%%%%%%%%%%%%%%%%%%%%%%%%%%%%%%%%%%%%%%%%%%%%%%%%%%%%%%%%%%%%%%%%%%%%%
%%%%%%%%%%%%%%%%%%%% dedicated to boxed/framed environements %%%%%%%%%%%%%%
%%%%%%%%%%%%%%%%%%%%%%%%%%%%%%%%%%%%%%%%%%%%%%%%%%%%%%%%%%%%%%%%%%%%%%%%%%%
\newtheoremstyle{bluenumbox}% % Theorem style name
{0pt}% Space above
{0pt}% Space below
{\normalfont}% % Body font
{}% Indent amount
{\small\bf\sffamily\color{blue}}% % Theorem head font
{\;}% Punctuation after theorem head
{0.25em}% Space after theorem head
{\small\sffamily\color{blue}\thmname{#1}\nobreakspace\thmnumber{\@ifnotempty{#1}{}\@upn{#2}}% Theorem text (e.g. Theorem 2.1)
\thmnote{\nobreakspace\the\thm@notefont\sffamily\bfseries\color{black}---\nobreakspace#3.}} % Optional theorem note
\renewcommand{\qedsymbol}{$\blacksquare$}% Optional qed square

\newtheoremstyle{blacknumex}% Theorem style name
{5pt}% Space above
{5pt}% Space below
{\normalfont}% Body font
{} % Indent amount
{\small\bf\sffamily}% Theorem head font
{\;}% Punctuation after theorem head
{0.25em}% Space after theorem head
{\small\sffamily{\tiny\ensuremath{\blacksquare}}\nobreakspace\thmname{#1}\nobreakspace\thmnumber{\@ifnotempty{#1}{}\@upn{#2}}% Theorem text (e.g. Theorem 2.1)
\thmnote{\nobreakspace\the\thm@notefont\sffamily\bfseries---\nobreakspace#3.}}% Optional theorem note

\newtheoremstyle{blacknumbox} % Theorem style name
{0pt}% Space above
{0pt}% Space below
{\normalfont}% Body font
{}% Indent amount
{\small\bf\sffamily}% Theorem head font
{\;}% Punctuation after theorem head
{0.25em}% Space after theorem head
{\small\sffamily\thmname{#1}\nobreakspace\thmnumber{\@ifnotempty{#1}{}\@upn{#2}}% Theorem text (e.g. Theorem 2.1)
\thmnote{\nobreakspace\the\thm@notefont\sffamily\bfseries---\nobreakspace#3.}}% Optional theorem note

%%%%%%%%%%%%%%%%%%%%%%%%%%%%%%%%%%%%%%%%%%%%%%%%%%%%%%%%%%%%%%%%%%%%%%%%%%%
%%%%%%%%%%%%% dedicated to non-boxed/non-framed environements %%%%%%%%%%%%%
%%%%%%%%%%%%%%%%%%%%%%%%%%%%%%%%%%%%%%%%%%%%%%%%%%%%%%%%%%%%%%%%%%%%%%%%%%%
\newtheoremstyle{bluenum}% % Theorem style name
{5pt}% Space above
{5pt}% Space below
{\normalfont}% % Body font
{}% Indent amount
{\small\bf\sffamily\color{blue}}% % Theorem head font
{\;}% Punctuation after theorem head
{0.25em}% Space after theorem head
{\small\sffamily\color{blue}\thmname{#1}\nobreakspace\thmnumber{\@ifnotempty{#1}{}\@upn{#2}}% Theorem text (e.g. Theorem 2.1)
\thmnote{\nobreakspace\the\thm@notefont\sffamily\bfseries\color{black}---\nobreakspace#3.}} % Optional theorem note
\renewcommand{\qedsymbol}{$\blacksquare$}% Optional qed square
\makeatother

% Defines the theorem text style for each type of theorem to one of the three styles above
\newcounter{dummy} 
\numberwithin{dummy}{section}
\theoremstyle{bluenumbox}
\newtheorem{theoremeT}[dummy]{Theorem}
\newtheorem{problem}{Problem}[chapter]
\newtheorem{exerciseT}{Exercise}[chapter]
\theoremstyle{blacknumex}
\newtheorem{exampleT}{Example}[chapter]
\theoremstyle{blacknumbox}
\newtheorem{vocabulary}{Vocabulary}[chapter]
\newtheorem{definitionT}{Definition}[section]
\newtheorem{corollaryT}[dummy]{Corollary}
\theoremstyle{bluenum}
\newtheorem{proposition}[dummy]{Proposition}

%----------------------------------------------------------------------------------------
%	DEFINITION OF COLORED BOXES
%----------------------------------------------------------------------------------------

\RequirePackage[framemethod=default]{mdframed} % Required for creating the theorem, definition, exercise and corollary boxes

% Theorem box
\newmdenv[skipabove=7pt,
skipbelow=7pt,
backgroundcolor=black!5,
linecolor=blue,
innerleftmargin=5pt,
innerrightmargin=5pt,
innertopmargin=5pt,
leftmargin=0cm,
rightmargin=0cm,
innerbottommargin=5pt]{tBox}

% Exercise box	  
\newmdenv[skipabove=7pt,
skipbelow=7pt,
rightline=false,
leftline=true,
topline=false,
bottomline=false,
backgroundcolor=blue!10,
linecolor=blue,
innerleftmargin=5pt,
innerrightmargin=5pt,
innertopmargin=5pt,
innerbottommargin=5pt,
leftmargin=0cm,
rightmargin=0cm,
linewidth=4pt]{eBox}	

% Definition box
\newmdenv[skipabove=7pt,
skipbelow=7pt,
rightline=false,
leftline=true,
topline=false,
bottomline=false,
linecolor=blue,
innerleftmargin=5pt,
innerrightmargin=5pt,
innertopmargin=0pt,
leftmargin=0cm,
rightmargin=0cm,
linewidth=4pt,
innerbottommargin=0pt]{dBox}	

% Corollary box
\newmdenv[skipabove=7pt,
skipbelow=7pt,
rightline=false,
leftline=true,
topline=false,
bottomline=false,
linecolor=gray,
backgroundcolor=black!5,
innerleftmargin=5pt,
innerrightmargin=5pt,
innertopmargin=5pt,
leftmargin=0cm,
rightmargin=0cm,
linewidth=4pt,
innerbottommargin=5pt]{cBox}

% Creates an environment for each type of theorem and assigns it a theorem text style from the "Theorem Styles" section above and a colored box from above
\newenvironment{theorem}{\begin{tBox}\begin{theoremeT}}{\end{theoremeT}\end{tBox}}
\newenvironment{exercise}{\begin{eBox}\begin{exerciseT}}{\hfill{\color{blue}\tiny\ensuremath{\blacksquare}}\end{exerciseT}\end{eBox}}				  
\newenvironment{definition}{\begin{dBox}\begin{definitionT}}{\end{definitionT}\end{dBox}}	
\newenvironment{example}{\begin{exampleT}}{\hfill{\tiny\ensuremath{\blacksquare}}\end{exampleT}}		
\newenvironment{corollary}{\begin{cBox}\begin{corollaryT}}{\end{corollaryT}\end{cBox}}	

%----------------------------------------------------------------------------------------
%	REMARK ENVIRONMENT
%----------------------------------------------------------------------------------------

\newenvironment{remark}{\par\vspace{10pt}\small % Vertical white space above the remark and smaller font size
\begin{list}{}{
\leftmargin=35pt % Indentation on the left
\rightmargin=25pt}\item\ignorespaces % Indentation on the right
\makebox[-2.5pt]{\begin{tikzpicture}[overlay]
\node[draw=blue!60,line width=1pt,circle,fill=blue!25,font=\sffamily\bfseries,inner sep=2pt,outer sep=0pt] at (-15pt,0pt){\textcolor{blue}{R}};\end{tikzpicture}} % Orange R in a circle
\advance\baselineskip -1pt}{\end{list}\vskip5pt} % Tighter line spacing and white space after remark

%----------------------------------------------------------------------------------------
%	SECTION NUMBERING IN THE MARGIN
%----------------------------------------------------------------------------------------

\makeatletter
\renewcommand{\@seccntformat}[1]{\llap{\textcolor{blue}{\csname the#1\endcsname}\hspace{1em}}}                    
\renewcommand{\section}{\@startsection{section}{1}{\z@}
{-4ex \@plus -1ex \@minus -.4ex}
{1ex \@plus.2ex }
{\normalfont\large\sffamily\bfseries}}
\renewcommand{\subsection}{\@startsection {subsection}{2}{\z@}
{-3ex \@plus -0.1ex \@minus -.4ex}
{0.5ex \@plus.2ex }
{\normalfont\sffamily\bfseries}}
\renewcommand{\subsubsection}{\@startsection {subsubsection}{3}{\z@}
{-2ex \@plus -0.1ex \@minus -.2ex}
{.2ex \@plus.2ex }
{\normalfont\small\sffamily\bfseries}}                        
\renewcommand\paragraph{\@startsection{paragraph}{4}{\z@}
{-2ex \@plus-.2ex \@minus .2ex}
{.1ex}
{\normalfont\small\sffamily\bfseries}}

%----------------------------------------------------------------------------------------
%	HYPERLINKS IN THE DOCUMENTS
%----------------------------------------------------------------------------------------

% For an unclear reason, the package should be loaded now and not later
\usepackage{hyperref}
\hypersetup{hidelinks,backref=true,pagebackref=true,hyperindex=true,colorlinks=false,breaklinks=true,urlcolor= blue,bookmarks=true,bookmarksopen=false,pdftitle={Title},pdfauthor={Author}}

%----------------------------------------------------------------------------------------
%	CHAPTER HEADINGS
%----------------------------------------------------------------------------------------

% The set-up below should be (sadly) manually adapted to the overall margin page septup controlled by the geometry package loaded in the main.tex document. It is possible to implement below the dimensions used in the goemetry package (top,bottom,left,right)... TO BE DONE

\newcommand{\thechapterimage}{}
\newcommand{\chapterimage}[1]{\renewcommand{\thechapterimage}{#1}}

% Numbered chapters with mini tableofcontents
\def\thechapter{\arabic{chapter}}
\def\@makechapterhead#1{
\thispagestyle{empty}
{\centering \normalfont\sffamily
\ifnum \c@secnumdepth >\m@ne
\if@mainmatter
\startcontents
\begin{tikzpicture}[remember picture,overlay]
\node at (current page.north west)
{\begin{tikzpicture}[remember picture,overlay]
\node[rounded corners=5pt,anchor=north west,inner sep=0pt] at (0,0) {\includegraphics[width=\paperwidth]{\thechapterimage}};
%%%%%%%%%%%%%%%%%%%%%%%%%%%%%%%%%%%%%%%%%%%%%%%%%%%%%%%%%%%%%%%%%%%%%%%%%%%%%%%%%%%%%
%Tamanho do mini table of contents
\fill[color=blue!10!white,opacity=.6] (.5cm,0) rectangle (9.9cm,-13cm);
%Conteúdo do mini table of contents 
\node[rounded corners=10pt,anchor=north west] at (1.1cm,.35cm) {\parbox[t][10cm][t]{10cm}{\huge\bfseries\flushleft \printcontents{l}{1}{\setcounter{tocdepth}{2}}}};
%Titulo do capitulo
\draw[anchor=west] (10cm,-10cm) node [rounded corners=5pt,fill=blue!10!white,text opacity=1,draw=blue,draw opacity=1,line width=1.5pt,fill opacity=.6,inner sep=12pt]{\huge\sffamily\bfseries\textcolor{black}{\thechapter. #1\strut\makebox[20cm]{}}};
%%%%%%%%%%%%%%%%%%%%%%%%%%%%%%%%%%%%%%%%%%%%%%%%%%%%%%%%%%%%%%%%%%%%%%%%%%%%%%%%%%%%%
\end{tikzpicture}};
\end{tikzpicture}}
\par\vspace*{230\p@}
\fi
\fi}

% Unnumbered chapters without mini tableofcontents (could be added though) 
\def\@makeschapterhead#1{
\thispagestyle{empty}
{\centering \normalfont\sffamily
\ifnum \c@secnumdepth >\m@ne
\if@mainmatter
\begin{tikzpicture}[remember picture,overlay]
\node at (current page.north west)
{\begin{tikzpicture}[remember picture,overlay]
\node[anchor=north west,inner sep=0pt] at (0,0) {\includegraphics[width=\paperwidth]{\thechapterimage}};
\draw[anchor=west] (2cm,-9cm) node [rounded corners=5pt,fill=blue!10!white,fill opacity=.6,inner sep=12pt,text opacity=1,draw=blue,draw opacity=1,line width=1.5pt]{\huge\sffamily\bfseries\textcolor{black}{#1\strut\makebox[22cm]{}}};
\end{tikzpicture}};
\end{tikzpicture}}
\par\vspace*{230\p@}
\fi
\fi
}
\makeatother % Insert the commands.tex file which contains the majority of the structure behind the template

\begin{document}

\let\cleardoublepage\clearpage

\renewcommand{\chaptername}{Capítulo}
\renewcommand{\figurename}{Fig.}

%----------------------------------------------------------------------------------------
%	TITLE PAGE
%----------------------------------------------------------------------------------------
\begingroup
	\thispagestyle{empty}
	
	\AddToShipoutPicture*{\put(0,0){\includegraphics[scale=1]{capa}}} % Image background
	
	\AddToShipoutPicture*{\put(116,650){\includegraphics[scale=.75]{brasao.png}}} % Image background
	
	\AddToShipoutPicture*{\put(244,200){\includegraphics[scale=0.2]{ifgvertical}}} % Image background
	
	\vspace*{4.5cm}
	
	\centering
	\par
	{\Huge Projeto Pedagógico}\vspace*{1.5cm}
	\par
	\fontsize{40}{40}
	\selectfont
	Tecnologia em Análise e Desenvolvimento de Sistemas
	\vspace*{10cm}
	\par
	{\Huge 2020}
	\par
\endgroup
\pagebreak

%----------------------------------------------------------------------------------------
%	PEOPLE PAGE
%----------------------------------------------------------------------------------------
\chapterimage{banner3} % Chapter heading image
\begin{center}
	\par
	{\large PRESIDENTE DA REPÚBLICA \\ Jair Messias Bolsonaro}\vspace*{1cm}
	\par
	{\large MINISTRO DA EDUCAÇÃO \\ Nome do Ministro}\vspace*{1cm}
	\par
	{\large SECRETÁRIO DE EDUCAÇÃO PROFISSIONAL E TECNOLÓGICA \\ Alexandro Ferreira de Souza}\vspace*{1cm}
	\par
	{\large REITOR DO INSTITUTO FEDERAL DE GOIÁS \\ Jerônimo Rodrigues da Silva}\vspace*{1cm}
	\par
	{\large PRÓ-REITOR DE PESQUISA E PÓS-GRADUAÇÃO \\ Ruberley Rodrigues de Souza}\vspace*{1cm}
	\par
	{\large DIRETORIA DE PÓS-GRADUAÇÃO \\ Clarinda Aparecida da Silva}\vspace*{1cm}
	\par
	{\large COORDENADOR DO CURSO \\ Nome do Coordenador}\vspace*{1cm}
\end{center}

\chapterimage{banner3} % Chapter heading image
\renewcommand\contentsname{Sumário}
\tableofcontents

%----------------------------------------------------------------------------------------
%	CHAPTER
%----------------------------------------------------------------------------------------
\chapterimage{01.jpg} % Chapter heading image
\chapter{Apresentação}\label{apresentacao}
\vspace{6em}
\begin{flushright}
	\textit{\textcolor{white}{Um bonita citação...}}
\end{flushright}
\vspace{12em}

\todo[inline]{NNNNNNNNNNNNN}~\parencite{Resolucao3De2002}


%-----------------------------------------------
\newpage  
%------------------------------------------------
\section{Identificação do Curso}
\begin{itemize}
	\item \textbf{Instituição Ofertante:} Instituto Federal de Educação, Ciência e Tecnologia de Goiás
	\item \textbf{Nome do curso:} Tecnologia em Análise e Desenvolvimento de Sistemas
	\item \textbf{Carga Horária do Curso:} 2425 horas
	\item \textbf{Forma de oferta:} Presencial
	\item \textbf{Duração:} 2,5 anos
	\item \textbf{Número de Vagas:} 30 vagas anuais
	\item \textbf{Local de Oferta:} Instituto Federal de Goiás - Câmpus Formosa
	\item \textbf{Reitor:} NNNNNNNNNNNNN
	\item \textbf{Pró-Reitora de Ensino:} NNNNNNNNNNNNN
	\item \textbf{Pró-Reitor de Pesquisa, Pós-Graduação e Inovação:} NNNNNNNNNNNNN
	\item \textbf{Diretoria de Pós-Graduação:} NNNNNNNNNNNNN
\end{itemize}

\section{Elaboração do Projeto de Curso}
\begin{itemize}[label=\bfseries]
	\item ...
	\item Waldeyr Mendes Cordeiro da Silva
\end{itemize}

%----------------------------------------------------------------------------------------
%	CHAPTER
%----------------------------------------------------------------------------------------
\chapterimage{02.jpg} % Chapter heading image
\chapter{Introdução}\label{introducao}
\vspace{6em}
\begin{flushright}
	\textit{\textcolor{white}{Um bonita citação...}}
\end{flushright}
\vspace{12em}

\todo[inline]{NNNNNNNNNNNNN}

\section{Justificativa}

\todo[inline]{NNNNNNNNNNNNN}

\section{Público Alvo}

\todo[inline]{NNNNNNNNNNNNN}

\section{Objetivos}\label{objetivos}

\subsection{Objetivo Geral}

\todo[inline]{NNNNNNNNNNNNN}

\subsection{Objetivos Específicos}

\begin{itemize}
\item NNNNNNNNNNNNN 
\end{itemize}


\section{Perfil do Egresso}
O Tecnólogo em Análise e Desenvolvimento de Sistemas:
\begin{enumerate}
	\item Analisa, projeta, desenvolve, testa, implanta e mantém sistemas computacionais de informação. 
	\item Avalia, seleciona, especifica e utiliza metodologias, tecnologias e ferramentas da Engenharia de Software, linguagens de programação e bancos de dados. 
	\item Coordena equipes de produção de softwares.
	\item Vistoria, realiza perícia, avalia, emite laudo e parecer técnico em sua área de formação.
	\item Compreende e domina os fundamentos científico-tecnológicos do seu trabalho.
	\item Respeita as diversidades e os direitos humanos, tem atitude ética e responsabilidade social e ambiental no trabalho e fora dele.
	\item Compreende a importância e impacto do seu trabalho.
	\item Usa a linguagem para a cidadania e profissão.
	\item Trabalha em equipe e conhece planejamento estratégico, iniciativa e liderança.
\item 	Atualiza-se mantendo-se criativo e responsável.
\end{enumerate}



%----------------------------------------------------------------------------------------
%	CHAPTER
%----------------------------------------------------------------------------------------
\chapterimage{03.jpg} % Chapter heading image
\chapter{Organização do Curso}\label{organizacao}
\vspace{6em}
\begin{flushright}
	\textit{\textcolor{white}{Um bonita citação...}}
\end{flushright}
\vspace{12em}

\section{Requisitos para Acesso ao Curso}

\todo[inline]{NNNNNNNNNNNNN}

\section{Forma de Oferta}\label{carga}


\section{Matriz Curricular}\label{matriz}

\todo[inline]{NNNNNNNNNNNNN}

\begin{table}[]
	\centering
	\caption{Matriz Curricular do Curso Superior em Tecnologia em Análise e Desenvolvimento de Sistemas.}
		\label{tab:matriz}
		\resizebox{\textwidth}{!}{%
			\begin{tabular}{|l|l|c|c|c|c|c|c|}
				\hline
				\textbf{Núcleo}                               
				& \textbf{Disciplinas}                          & \textbf{1º semestre} & \textbf{2º semestre} & \textbf{3º semestre} & \textbf{4º semestre} & \textbf{5º semestre} & \textbf{CH}   \\ \hline
				\multirow{22}{*}{\textbf{Profissionalizante}} 
				& \nameref{disc:algoritmos}                     & 54                   &                      &                      &                      &                      & 54            \\ \cline{2-8} 
				& Fundamentos da Computação                     & 54                   &                      &                      &                      &                      & 54            \\ \cline{2-8} 
				& Engenharia de Software                        & 54                   &                      &                      &                      &                      & 54            \\ \cline{2-8} 
				& Sistemas Operacionais                         &                      & 54                   &                      &                      &                      & 54            \\ \cline{2-8} 
				& Bancos de Dados                               &                      & 54                   &                      &                      &                      & 54            \\ \cline{2-8} 
				& Bancos de Dados II                            &                      &                      & 54                   &                      &                      & 54            \\ \cline{2-8} 
				& Estruturas de Dados                           &                      & 54                   &                      &                      &                      & 54            \\ \cline{2-8} 
				& Programação para Web I                        &                      & 54                   &                      &                      &                      & 54            \\ \cline{2-8} 
				& Programação para Web II                       &                      &                      & 54                   &                      &                      & 54            \\ \cline{2-8} 
				& Análise e desenvolvimento de sistemas I       &                      &                      & 81                   &                      &                      & 81            \\ \cline{2-8} 
				& Análise e desenvolvimento de sistemas II      &                      &                      &                      & 81                   &                      & 81            \\ \cline{2-8} 
				& Análise e desenvolvimento de sistemas III     &                      &                      &                      &                      & 81                   & 81            \\ \cline{2-8} 
				& Redes de Computadores                         &                      &                      & 54                   &                      &                      & 54            \\ \cline{2-8} 
				& Administração de Serviços para Internet       &                      &                      & 54                   &                      &                      & 54            \\ \cline{2-8} 
				& Gerência e Governança em Tecnologia da Informação                    &                      &                      & 54                   &                      &                      & 54            \\ \cline{2-8} 
				& Verificação, Validação e Teste de Software    &                      &                      &                      & 27                   &                      & 27            \\ \cline{2-8} 				
				& Segurança da Informação                       &                      &                      &                      &                      & 54                   & 54            \\ \cline{2-8} 
				& Inteligência Artificial                       &                      &                      &                      &                      & 54                   & 54            \\ \cline{2-8} 
				& Componente Curricular Eletivo Técnico I       &                      & 81                   &                      &                      &                      & 81            \\ \cline{2-8} 
				& Componente Curricular Eletivo Técnico II      &                      &                      &                      & 81                   &                      & 81            \\ \cline{2-8} 
				& Componente Curricular Eletivo Técnico III     &                      &                      &                      &                      & 81                   & 81            \\ \cline{2-8} 
				& \multicolumn{6}{l|}{\textbf{Carga Horária Disciplinas Técnicas}}                                                                                                 & \textbf{1269} \\ \hline
				\multirow{17}{*}{\textbf{Formação Integral}}  
				& Matemática                                    & 54                   &                      &                      &                      &                      & 54            \\ \cline{2-8} 
				& Inglês Acadêmico                              & 54                   &                      &                      &                      &                      & 54            \\ \cline{2-8} 
				& Empreendedorismo                              &                      &                      &                      & 27                   &                      & 27            \\ \cline{2-8} 				
				& Leitura e Produção de Textos                  & 54                   &                      &                      &                      &                      & 54            \\ \cline{2-8} 
				& Sociologia do Trabalho                        & 27                   &                      &                      &                      &                      & 27            \\ \cline{2-8} 
				& Educação Ambiental                            & 27                   &                      &                      &                      &                      & 27            \\ \cline{2-8} 
				& Legislação Aplicada                           & 27                   &                      &                      &                      &                      & 27            \\ \cline{2-8} 
				& Probabilidade e Estatística                   &                      & 54                   &                      &                      &                      & 54            \\ \cline{2-8} 
				& Componente Curricular Eletivo Politécnico I   &                      & 54                   &                      &                      &                      & 54            \\ \cline{2-8} 
				& Componente Curricular Eletivo Politécnico II  &                      &                      & 54                   &                      &                      & 54            \\ \cline{2-8} 
				& Componente Curricular Eletivo Politécnico III &                      &                      & 54                   &                      &                      & 54            \\ \cline{2-8} 
				& Componente Curricular Eletivo Politécnico IV  &                      &                      &                      & 54                   &                      & 54            \\ \cline{2-8} 
				& Componente Curricular Eletivo Politécnico V   &                      &                      &                      &                      & 54                   & 54            \\ \cline{2-8} 
				& Componente Curricular Eletivo Politécnico VI  &                      &                      &                      &                      & 54                   & 54            \\ \cline{2-8} 
				& Componente Curricular Eletivo Humanas I       &                      &                      &                      & 27                   &                      & 27            \\ \cline{2-8} 
				& Componente Curricular Eletivo Humanas II      &                      &                      &                      &                      & 27                   & 27            \\ \cline{2-8} 
				& Metodologia e Iniciação Científica            &                      &                      &                      & 54                   &                      & 54            \\ \cline{2-8} 
				& \multicolumn{6}{l|}{\textbf{Carga Horária Disciplinas de Formação Integral}}                                                                                     & \textbf{756}  \\ \hline
				\multicolumn{7}{|l|}{\textbf{Carga Horária Disciplinas}}                                                                                                           & \textbf{2025} \\ \hline
				\multirow{2}{*}{\textbf{TCC}}                 
				& TCC I                                         &                      &                      &                      & 54                   &                      & -             \\ \cline{2-8} 
				& TCC II                                        &                      &                      &                      &                      & 54                   & -             \\ \hline
				\multicolumn{7}{|l|}{\textbf{Atividades complementares}}                                                                                                                                                         & \textbf{400}  \\ \hline
				\multicolumn{7}{|l|}{\textbf{Carga Horária Total}}                                                                                                                                                               & \textbf{2425} \\ \hline
			\end{tabular}%
		}
	\end{table}




\section{Metodologia de Ensino-Aprendizagem}\label{metodologia}


\subsection{Trabalhos Discentes}\label{trabdiscentes}

\subsubsection{Disciplinas}\label{dsiciplinas}

\todo[inline]{NNNNNNNNNNNNN}

\subsubsection{Atividades Não Presenciais}

Definição. Arcabouço jurídico.


\subsubsection{Avaliação}

\todo[inline]{NNNNNNNNNNNNN}

\section{Certificação}

O Certificado será emitido pelo Instituto Federal de Educação, Ciência e Tecnologia de Goiás, nos termos da Resolução CNE/CES n.º 1, de 8 de junho de 2007.	
Para obter o Certificado de graduado em ``Tecnologia em Análise e Desenvolvimento de Sistemas'', o discente deverá satisfazer as seguintes exigências:
\begin{itemize}
	\item Ser aprovado em todas as disciplinas do curso com nota mínima igual a 6,0 (seis) e freqüência igual ou superior a 75\% da carga horária da disciplina;
	\item Ser aprovado em defesa pública do trabalho de conclusão de curso perante uma banca composta por, no mínimo, três professores (orientador, mais dois professores convidados);
	\item Possuir pelo menos um certificado que comprove a apresentação (pôster ou oral) de resultados relacionados ao trabalho de conclusão de curso em evento científico (congressos, seminários, simpósios);
	\item Quitação de todas as obrigações junto ao Câmpus Formosa do Instituto Federal de Educação, Ciência e Tecnologia de Goiás;
\end{itemize}

%----------------------------------------------------------------------------------------
%	CHAPTER
%----------------------------------------------------------------------------------------
%------------------------------------------------
\chapterimage{04.jpg} % Chapter heading image
\chapter{Ementas}\label{ementas}
\vspace{6em}
\begin{flushright}
	\textit{\textcolor{white}{Um bonita citação...}}
\end{flushright}
\vspace{12em}

\todo[inline]{NNNNNNNNNNNNN}

\newpage
\section{Algoritmos}\label{disc:algoritmos}

\begin{itemize}
	\item \textbf{Carga horária (hora/aula):} 54
	\item \textbf{Docente Responsável:}~\nameref{WaldeyrMendes}
	\item \textbf{Ementa:} 
	Conceitos de algoritmos;
	Conceitos de linguagens de programação;
	Constantes e variáveis;
	Tipos de dados;
	Operadores e expressões aritméticas, lógicas e literais; 
	Comandos básicos;
	Estruturas condicionais e de repetição;
	Vetores e matrizes;
	Estruturas de dados básicas;
	Modularização;
	Recursividade;
	Algoritmos e meio ambiente;
	\item \textbf{Bibliografia básica}
	\begin{enumerate}
		\item CORMEN T. H., LEISERSON C. E., RIVEST R. L.,STEIN C..; Algoritmos: Teoria e prática. Editora Câmpus, Tradução da 3a ed. americana, São Paulo-SP, 2012.
		\item 
	\end{enumerate}
	\item \textbf{Bibliografia complementar}
	\begin{enumerate}
		\item NNNNNNNNNNNNN
	\end{enumerate}	
\end{itemize}



%----------------------------------------------------------------------------------------
%	CHAPTER X
%----------------------------------------------------------------------------------------
%------------------------------------------------
\chapterimage{05.jpg} % Chapter heading image
\chapter{Estrutura Física}\label{estrutura}
\vspace{6em}
\begin{flushright}
	\textit{\textcolor{white}{Um bonita citação...}}
\end{flushright}
\vspace{12em}

\section{Laboratórios de Informática}
Dois laboratórios de informática com capacidade para até 30 estudantes, com acesso à Internet, computadores com sistema operacional Linux, softwares diversos.


\section{Laboratório de Fisiologia Vegetal}
Equipado com: estufa de secagem, 3 estereoscópios, 3 microscópicos, geladeira, bancadas, 28 cadeiras, quadro e acervo didático (frutos, sementes e folhas herborizadas). 

\section{Laboratório de Bioqímica}
Equipado com: Balanças analítica e semi-analítica, chapas de aquecimento (com agitação magnética), analisador bioquímico, capela de fluxo laminar, agitadores de tubo de ensaio, banho-maria, bomba de vácuo, autoclave, estufas, destilador e deionizador de água e outros.

\section{Laboratório de Anatomia e Zoologia}
Equipado com: Bonecos anatômicos (de abdome) completos, conjuntos anatômicos artificiais de sistemas reprodutores femininos e masculinos, esqueletos completos (artificiais), amostras de animais (do cerrado e de outros biomas) conservados em frascos para visualização, animais empalhados, algumas peças anatômicas naturais de animais, lupas, microscópios, material para coleta de animais e saídas de campo, materiais e reagentes para o empalhamento de animais e outros.

\section{Laboratório de Microscopia e Microbiologia}
Equipado com:  25 microscópios e material para produção de lâminas (lâminas de corte, lâminas e lamínulas de vidro, corantes, fixadores, etc); Lupas, coleções de laminários e outros.

\section{Laboratório de Físico-Química}
Equipado com: pHmetros, destilador, capela de exaustão, estufa, banho-maria, balanças analítica e semi-analítica, deionizador, reator, aparelho de ponto de fusão,  e outros.

\section{Laboratório de Águas Residuais}
Equipado com: Condutivímetros, muflas, banho - maria, bomba de vácuo, analisador de oxigênio dissolvido, turbdímetro, estufa, balança, phmetro, destilador e outros.

\section{Laboratório de Ensino}
Espaço acadêmico voltado ao desenvolvimento e disseminação de tecnologias educacionais voltadas ao ensino de Ciências e Biologia.  Equipado com: acervo didático constituído por jogos, maquetes e representações físicas de organismos e processos biológicos.

\section{Laboratório de Física e Matemática}
O Laboratório de Física possui diversos equipamentos que contribui para o desenvolvimento das atividades experimentais nas áreas de mecânica, óptica, hidrostática, termologia e eletricidade.

\section{Biblioteca}
Biblioteca equipada com áreas de estudo individual e coletivo, 6 computadores com acesso ao portal de periódicos e acervo cerca de 7 mil exemplares, entre livros, livros em braile, cds, dvds e mapas;

\section{Teatro}
Teatro equipado com som e iluminação específica e acomodações para 320 pessoas sentadas;

\section{Outros Espaços}
3 salas para estudos coletivos e reuniões equipadas com mesas, cadeiras e televisor.


%----------------------------------------------------------------------------------------
%	CHAPTER X
%----------------------------------------------------------------------------------------
%------------------------------------------------
\chapterimage{06.jpg} % Chapter heading image
\chapter{Corpo Docente}\label{docentes}
\vspace{6em}
\begin{flushright}
	\textit{\textcolor{white}{Um bonita citação...}}
\end{flushright}
\vspace{12em}

\section{Waldeyr Mendes Cordeiro da Silva}\label{WaldeyrMendes}
\begin{itemize}
	\item Formação Básica: Sistemas de Informação e Ciências Biológicas 
	\item Titulação Máxima: Doutor em Ciências Biológicas (Bioinformática)
	\item Regime de Trabalho: Dedicação Exclusiva
	\item \includegraphics[scale=.03]{Pictures/lattes}~\href{http://lattes.cnpq.br/2391349697609978}{Lattes: http://lattes.cnpq.br/2391349697609978}
	\item \includegraphics[scale=.15]{Pictures/orcid}~\href{https://orcid.org/0000-0002-8660-6331}{ORCID: https://orcid.org/0000-0002-8660-6331}
\end{itemize}

% ----------------------------------------------------------------------------------------
% 	BIBLIOGRAPHY
% ----------------------------------------------------------------------------------------
%----------------------------------------------------------------------------------------
%	CHAPTER X
%----------------------------------------------------------------------------------------
%------------------------------------------------
\chapterimage{07.jpg} % Chapter heading image
%\chapter*{Referências Bibliográficas}
%\bibliography{bibliography}
%\renewcommand\bibname{Referências Bibliográficas}

\chapter*{Referências Bibliográficas}\label{referencias}
\vspace{6em}
\begin{flushright}
	\textit{\textcolor{white}{Um bonita citação...}}
\end{flushright}
\vspace{12em}
%\addcontentsline{toc}{chapter}{\textcolor{verde}{Bibliography}}
%\section{Books}
%\addcontentsline{toc}{section}{Books}
%\printbibliography[heading=bibempty,type=book]
%\section{Articles}
%\addcontentsline{toc}{section}{Articles}
%\printbibliography[heading=bibempty,type=article]
\printbibliography[heading=bibempty]


%----------------------------------------------------------------------------------------

%----------------------------------------------------------------------------------------
%	INDEX
%----------------------------------------------------------------------------------------
%
%\cleardoublepage
%\phantomsection
%\setlength{\columnsep}{0.75cm}
%\addcontentsline{toc}{chapter}{\textcolor{verde}{Index}}
%\printindex


%----------------------------------------------------------------------------------------
%	CHAPTER X
%----------------------------------------------------------------------------------------
%------------------------------------------------
\chapterimage{08.jpg} % Chapter heading image
\chapter{Anexos}\label{conheca}
\vspace{6em}
\begin{flushright}
	\textit{\textcolor{white}{}}
\end{flushright}
\vspace{12em}

\newpage
\section{Compatibilidade com a Matriz Curricular Anterior}
\indent

Com a atualização da Matriz Curricular do Curso Superior em Tecnologia em Análise e Desenvolvimento de Sistemas, estudantes matriculados na matriz anterior poderão migrar para a nova matriz respeitadas as compatibilidades definidas na Tabela~\ref{tab:compatibilidade}.
\begin{table}[H]
	\centering
	\caption{Tabela de compatibilidade de disciplinas entre a matriz curricular antiga e a atualizada.}
	\label{tab:compatibilidade}
	\resizebox{\textwidth}{!}{%
		\begin{tabular}{|l|l|}
			\hline
			\textbf{Disciplinas da Matriz Atuaizada}  & \textbf{Disciplinas da Matriz antiga}                                                            \\ \hline
			Algoritmos                                & Algoritmos                                                                                       \\ \hline
			Fundamentos da Computação                 & Fundamentos da Computação                                                                        \\ \hline
			Engenharia de Software                    & Engenharia de Software                                                                           \\ \hline
			Matemática                                & Matemática Elementar                                                                             \\ \hline
			Inglês Acadêmico                          & -                                                                                                \\ \hline
			Leitura e Produção de Textos              & -                                                                                                \\ \hline
			Sociologia do Trabalho                    & Sociologia do Trabalho                                                                           \\ \hline
			Educação Ambiental                        & Educação Ambiental                                                                               \\ \hline
			Legislação Aplicada                       & Ética e Legislação Aplicada à Informática                                                        \\ \hline
			Sistemas Operacionais                     & Sistemas Operacionais                                                                            \\ \hline
			Bancos de Dados                           & Banco de Dados I                                                                                 \\ \hline
			Estruturas de Dados                       & Estrutura de Dados I                                                                             \\ \hline
			Programação para Web I                    & Programação para Web I                                                                           \\ \hline
			Componente Curricular Eletivo Técnico     & Métodos e Técnicas de Programação                                                                \\ \hline
			Componente Curricular Eletivo Politécnico & Lógica Computacional                                                                             \\ \hline
			Probabilidade e Estatística               & Introdução à Probabilidade e Estatística + Cálculo Diferencial e Integral                        \\ \hline
			Análise e desenvolvimento de sistemas I   & Programação Orientada a Objetos                                                                  \\ \hline
			Redes de Computadores                     & Redes de Computadores                                                                            \\ \hline
			Programação para Web II                   & Programação para Web II                                                                          \\ \hline
			Gerência de Projetos                      & Gerência de Projetos                                                                             \\ \hline
			Administração de Serviços para Internet   & Administração de Serviços para Internet                                                          \\ \hline
			Componente Curricular Eletivo Politécnico & Cálculo Diferencial e Integral                                                                   \\ \hline
			Componente Curricular Eletivo Politécnico & Banco de Dados II                                                                                \\ \hline
			Componente Curricular Eletivo Humanas     & Optativa 1                                                                                       \\ \hline
			Análise e desenvolvimento de sistemas II  & Arquitetura e Projeto de Software                                                                \\ \hline
			Programação para Dispositivos Móveis      & -                                                                                                \\ \hline
			Gerência e Governança em Tecnologia da Informação                & -                                                                                                \\ \hline
			Componente Curricular Eletivo Técnico     & Estrutura de Dados II                                                                            \\ \hline
			Metodologia e Iniciação Científica        & Metodologia da Pesquisa Científica + Tópicos Avançados 1                                         \\ \hline
			Componente Curricular Eletivo Politécnico & Engenharia de Requisitos                                                                         \\ \hline
			TCC I                                     & TCC I                                                                                            \\ \hline
			Análise e desenvolvimento de sistemas III & Análise Orientada a Objetos + Fundamentos de Sistemas de Informação + Interface Homem Computador \\ \hline
			Segurança da Informação                   & Segurança e Auditoria de Sistemas                                                                \\ \hline
			Inteligência Artificial                   & -                                                                                                \\ \hline
			Componente Curricular Eletivo Técnico     & Sistemas Distribuídos                                                                            \\ \hline
			Componente Curricular Eletivo Politécnico & Computação Gráfica e Sistemas Multimídia                                                         \\ \hline
			Componente Curricular Eletivo Politécnico & Qualidade de Software                                                                            \\ \hline
			Componente Curricular Eletivo Humanas     & Optativa 2                                                                                       \\ \hline
			TCC II                                    & TCC II                                                                                           \\ \hline
		\end{tabular}%
	}
\end{table}






\end{document}